\section{Low SNR Mono-pulse Feasibility}
\todo{This is an opportunity to bring in really established papers on monopulse radar and explain what they do to maximise SNR. It may be hard to find gaps that aren't obvious.}
\todo{Develop the below further}

In addition to a lower SNR, a key limitation of SWaP-constrained monopulse radars is that a reduced aperture degrades angular resolution, increasing the likelihood that multiple targets occupy the same angular cell and that angle estimates begin reflecting the centroid of the target group. An et al. \cite{an2024_time_varying_angle_estimation_unresolved_extended_targets_monopulse} propose an efficient method for resolving two unresolved targets in a monopulse cell, but evaluation is limited to simulated white-noise conditions and does not address integration into an end-to-end range–Doppler detection. The method also has explicit failure modes (same azimuth/elevation; three or more targets in-cell). To address these gaps, our research will test this algorithm with real simulation data, where noise is not necessarily white, and compare how the incorporation of the algorithm into a tracking system improves track quality. 

Complementing this, Aubry et al. \cite{aubry2020_single_pulse_simultaneous_detection_angle_estimation_multichannel} demonstrate that detection and angle estimation can be treated jointly in a multichannel array, rather than as a strict two-stage pipeline (detect then estimate), which is attractive under SWaP conditions where false alarms and low SNR make early thresholding brittle. The research demonstrates an algorithm that fuses data between channels before angle tracking occurs, increasing sensitivity to targets that are not visible in any single channel. The research exclusively uses single simulated targets in Gaussian interference to validate the algorithm, leaving a gap in characterising the performance under real data. The proposed project will address these gaps by applying the algorithm to alternate noise regimes and multiple target scenarios. 
